\section{Related Work}
\label{sec:related}

\subsection{Old Film and Video Restoration}
\label{subsec:ofr}

Old film restoration differs from general video restoration in that it must jointly address structured artifacts (scratches, dust, blotches), photometric degradation (flicker, grain, fading), and quality degeneration (blur, noise). Traditional approaches decompose the problem into detection and inpainting stages, relying on hand-crafted features and interpolation-based filling~\cite{kokaram2004missing} or sparse collaborative filtering~\cite{dabov2007bm3d,he2013guided}. While computationally efficient, these methods lack semantic understanding of video content and leave photometric degradations unaddressed. Deep learning shifted old film restoration toward end-to-end frameworks: DeepRemaster~\cite{iizuka2019deepremaster} processes spatiotemporal volumes with 3D convolutions, RTN~\cite{wan2022rtn} introduces bidirectional recurrent propagation with a learnable soft defect mask, and RRTN~\cite{lin2024rrtn} strengthens alignment via recursive transformer networks with second-order grid propagation. MambaOFR~\cite{mao2025mambaofr} replaces computationally expensive blocks with MambaIR-style~\cite{guo2024mambair} state-space modules and establishes the SRWOV benchmark for dedicated old film restoration evaluation.

General video restoration methods such as EDVR~\cite{wang2019edvr}, BasicVSR~\cite{chan2021basicvsr}, BasicVSR++~\cite{chan2022basicvsrpp}, VRT~\cite{liang2024vrt}, and RVRT~\cite{liang2022rvrt} provide strong temporal modeling, but they do not incorporate defect handling specific to old film restoration. Recent diffusion and foundation-prior restoration systems, including DiffBIR~\cite{lin2024diffbir}, SUPIR~\cite{yu2024supir}, and SeeSR~\cite{wu2024seesr}, improve perceptual detail for blind image restoration or real-world super-resolution, but they are not designed around temporal old-film defects or SRWOV-style video restoration protocols. Despite this progress, existing old film restoration methods still encode structural completion, regional repair, and fine-detail suppression through a single propagated latent space. This motivates our multi-granularity decomposition approach.

\subsection{State Space Models for Vision}
\label{subsec:ssm}

The computational cost of self-attention scales quadratically with sequence length, posing a practical barrier for high-resolution video processing. Structured state space sequence models~\cite{gu2022s4} reformulate sequence modeling through continuous-time state transitions, reducing complexity from $O(L^2)$ to $O(L)$ while maintaining a global receptive field. Mamba~\cite{gu2024mamba} further introduces selective state transitions with hardware-aware recurrence, enabling content-dependent filtering. The adaptation of SSMs to visual domains has produced strong baselines: MambaIR~\cite{guo2024mambair} demonstrates 2D selective scanning for image restoration, VMamba~\cite{liu2024vmamba} builds hierarchical feature pyramids with four scanning routes, and VSRM~\cite{tran2025vsrm} alternates spatial and temporal Mamba blocks for video super-resolution at linear cost.

These properties are well-suited for old film restoration, which couples long temporal dependency, large spatial damage, and high-resolution input. MambaOFR~\cite{mao2025mambaofr} validates that Mamba-based propagation outperforms transformer counterparts on both efficiency and quality. However, linear-cost long-range aggregation alone does not decide what restoration evidence each feature position retains; the backbone treats all degradation scales uniformly. Our work augments the SSM feature flow with explicit multi-granularity descriptors rather than replacing the backbone.

\subsection{Multi-Scale and Multi-Branch Representations}
\label{subsec:multigran}

A single feature scale often cannot encode both global structure and fine-grained local evidence simultaneously. Feature Pyramid Networks~\cite{lin2017fpn} and HRNet~\cite{wang2021hrnet} establish that explicit multi-scale decomposition improves dense prediction, while MIRNet-v2~\cite{zamir2022mirnetv2} shows that complementary multi-scale context enhances low-level reconstruction quality. In person re-identification, PCB~\cite{sun2018pcb} decomposes features along spatial stripes, MGN~\cite{wang2018mgn} combines global and local branches at different partition granularities, and ReIDMamba~\cite{gu2025reidmamba} adapts this paradigm to state space models. These methods show that branch structure can expose complementary evidence, but their supervision and feature semantics target recognition rather than restoration.

\subsection{Diversity and Ranking Regularization}
\label{subsec:diversity_ranking}

The consistent lesson across these works is that extra branches provide benefit only when they capture complementary cues rather than correlated copies of the same evidence. Contrastive learning~\cite{chen2020mocov2} provides a general diversity mechanism, but it does not directly regulate pairwise ordering behavior across branches: two branches can produce different feature magnitudes yet rank all candidates identically, yielding diversity in norm but redundancy in discriminative structure. Differentiable Kendall's Tau~\cite{zheng2023diffkendall} enables gradient-based penalization of correlated rankings, a property that naturally applies to flow-guided temporal correspondences in old film restoration.

\subsection{Representation Symmetry in Restoration}
\label{subsec:symmetry_ai}

In AI-based image processing, symmetry can be interpreted beyond visible object geometry. A learned representation may be symmetric when complementary evidence sources remain balanced, comparable, and consistently organized across samples; asymmetry appears when optimization favors one dominant cue. Feature-fusion and correlation-based representation studies show that multi-view or multi-branch features need explicit organization rather than simple concatenation~\cite{yu2021multiview_tensor_fusion,qian2020multiview_matrix_correlation}, and feature-wise modulation provides a practical way to inject conditioning signals into image-processing networks~\cite{perez2018film}. For old film restoration, this view is useful because film damage is inherently asymmetric across regions and frames. In MgfrOFR this representation symmetry is made precise as two group invariances of the training objective---branch-permutation invariance under the symmetric group $\mathcal{S}_G$ and direction-exchange invariance under the bidirectional $\mathbb{Z}_2$ action (Section~\ref{subsec:symmetry_property}).

Our design is inspired by this line of work, but the restoration setting changes the formulation in three important ways. First, the branch tokens are inserted into propagated restoration features rather than identity-oriented recognition embeddings. Second, the branch summaries are fed back into the restoration stream through channel-wise recalibration, so they affect reconstruction rather than only auxiliary supervision. Third, positives and negatives are defined by temporal correspondence within and across clips rather than person identities. The analysis above highlights two gaps in old film restoration: the absence of explicit multi-granularity feature decomposition within the temporal propagation pipeline, and the lack of diversity regularization to ensure that parallel feature branches capture complementary restoration evidence. Unlike multi-scale pyramids, token/stripe recognition models, and generic diversity losses, MgfrOFR is the only design considered here whose branch structure carries an explicit permutation-and-exchange symmetry rather than acting on feature magnitude alone. MgfrOFR addresses both gaps within a unified framework, as detailed in Section~\ref{sec:method}.
